\section{A proposed research program}
\label{research:sec:program}

The results suggest concrete next problems. They are ordered by their
direct dependence on the new structure, rather than by a claim that
every direction has equal difficulty. Statements phrased as questions
below are not theorems of this article.

\subsection{Cyclotomic coefficient systems}

\paragraph{1. Determine the residual kernel at general level.}
Theorem~\ref{fg:thm:universal} removes every same-point sector
canonically. What is the remaining Hilbert series of
$\ker A_{N,w}^{\neg g}$ for fixed $N$? At levels three and four,
linear substitutions reduce the problem to small reflection groups.
At levels five, six, and eight, the next step is to write the remaining
color-polynomial equations, decompose them under color automorphisms,
and look for a finite invariant-theoretic presentation. Exact matrices
are valuable for selecting a candidate Hilbert series, but a proof
must identify generators, relations, and all boundary defects.

\paragraph{2. Add distribution without losing the canonical boundary data.}
Which of the free sectors survive when distribution rows are added?
The universal lift is supported on three colors per root pair, which
makes the image of each distribution row explicitly computable.
A useful first result would be a closed formula for the rank increment
at prime-power levels. This would also clarify the relation between
the coefficient systems here and the broader standard-relation spaces
of \cite{Zhao2010}.

\paragraph{3. Construct short relation certificates from the explicit dual basis.}
The nullspace test decides whether a proposed row is derivable, but
does not by itself minimize the number or coefficient height of the
rows used in a proof. Can the polynomial normal form be upgraded to
a constructive reduction algorithm whose certificate length is
polynomial in the weight? The existing $S_4$ certificate shows why
the choice of extra argument transformations matters. A principled
comparison between short certificates and large eliminated systems
would be useful for future repository formalization.

\subsection{Formal and analytic Herglotz arithmetic}

\paragraph{4. Make every positive classification case constructive.}
The congruence test answers whether the invariant companion symbol
vanishes, but a Bloch-group descent argument does not necessarily
produce a compact five-term derivation. Find an explicit logarithmic
formula, with a controllable number of terms, for all even--odd pairs
obeying the two congruences. The families
$J(e/(e^2-1))$ and $J(e/(e^2+1))$ are natural first targets beyond
consecutive arguments. Their formal reductions are proved here;
their simplest analytic evaluations remain to be found.

\paragraph{5. Classify subfield descent for combinations of symbols.}
An individual nonzero $\beta_M(a)$ cannot descend to a proper real
subfield, but combinations can. Character decomposition of the Gram
image should classify which coefficient vectors yield an element of
$\Lambda^2A(E)$ for a prescribed real subfield $E$. A complete answer
would generalize the denominator-rank formula and could distinguish
the smallest field needed by an explicit reduction from the ambient
cyclotomic field.

\paragraph{6. Complete the rank problem for companion values at fixed denominator.}
Theorem~\ref{jnew:thm:denominator-rank} gives an exact denominator
obstruction rank and a corresponding lower bound for formal companion
values modulo rational dilogarithms. What additional rank is supplied
by the numerator-conductor blocks? A useful formulation should fix a
finite numerator set and a common cyclotomic field, then compute the
intersection of the norm kernels. It should not confuse independence
in a symbol space with numerical independence of the evaluated periods.

\paragraph{7. Replace the dyadic difference by other Hecke combinations.}
The proof uses a cyclic Laplacian for multiplication by $2$ and a
quadratic field norm between adjacent dyadic conductors. Analogous
operators for an odd prime should interact with the Hecke functional
equations of \cite{RadchenkoZagier} and the character-weighted integrals
of \cite{DixitSathyanarayanaSharan2024}. Identify a natural companion
whose formal reduction is governed by a similarly explicit group
operator, and determine its exact kernel. The first challenge is to
keep rational wedge terms and conductor cancellation visible.

\subsection{Lerch bifurcations and uniform zero counts}

\paragraph{8. Prove or disprove uniqueness of the two-zero branch minimum.}
The smaller $n=2,k=1$ branch has at least one interior minimum and
an endpoint slope tending to $+\infty$. Does it have exactly one
critical point? Corollary~\ref{lerchshape:cor:n2} proves that the
larger branch is strictly decreasing, and that the smaller branch
increases after its unique crossing of $a=1$. The remaining question
therefore concerns only its motion below $1$. Differentiating the implicit equation gives
$a'=-G_\rho/G_a$; hence the first question reduces to the number of
solutions of $G=G_\rho=0$ on the first branch. A proof would require
a new sign-variation statement for this pair of kernels. A visually
single turn in the diagnostic plot is not enough to settle it.

\begin{conjecture}[Unique stationary point of the smaller branch]
\label{research:conj:minimum}
For $n=2,k=1$, the derivative $a_1'(\rho)$ vanishes at exactly one
point of $(0,1)$.
\end{conjecture}
The proved endpoint signs and Corollary~\ref{lerchshape:cor:n2}
would then identify that point as the unique global minimum, below
$\rho_c$. The numerical curve motivates the conjecture; no global
exclusion of additional stationary points is asserted here.

\paragraph{9. Locate and classify the forced interior events at $n=3,4$.}
At $k=1$, the small-$\rho$ counts and the certified $\rho=1$ counts
differ by two. We prove that an interior multiple zero occurs.
Determine whether there is exactly one event, and whether it is a
nondegenerate saddle-node, by certifying
\[
G=G_a=0,\qquad G_{aa}\ne0,\qquad G_\rho\ne0.
\]
Interval Newton or a two-variable argument principle could locate a
candidate, but global exclusion of further events must be supplied
separately. The current theorem deliberately makes no unverified
nondegeneracy claim.

\paragraph{10. Determine the optimal uniform threshold $K_n$.}
The present result gives $K_2=1$. The elementary endpoint forces
$K_n\geq n-1$: when $k<n-1$, the zero at $a=1$ is multiple and
uniform saturation by $n$ simple zeros fails at $\rho=0$.
The natural sharp question is whether $K_n=n-1$ for every $n\geq2$.
This equality is a candidate to test, not a conclusion of the existing
large-$k$ theorem. Exact sign certificates at selected separating
points, extended uniformly in $\rho$ by coefficient-sign arguments,
could prove the first new cases. A counterexample would reveal where
an interior loss of simplicity first occurs after the elementary
endpoint has already split completely.

\paragraph{11. Make the small-parameter classification effective in both indices.}
The sign of $\Delta_{n,k}$ is decidable by exact rational intervals
for $\log2$, but the proof gives no uniform lower bound for its
magnitude and no numerical value of $\rho_{n,k}$. Can one estimate
the distances of the roots of $P_{n,k}$ from $-\log2$ when $n,k$
grow? Such an estimate, together with a quantitative Rouché argument,
would turn the local classification into an explicit parameter range.
The resulting sign diagram may have a useful asymptotic description
when $k/n$ tends to a fixed ratio.

\paragraph{12. Compute the full Abel expansion of the zero branches.}
The leading singular coefficient is independent of $a$, while later
coefficients involve the tail expansion of $f_{n,k}(a+m)$ and the
regular moments. Derive the complete expansion in powers of $\tau$
and $\log\tau$, including collision terms from implicit inversion.
A particularly useful outcome would be a renormalized branch whose
successive singular terms can be subtracted algorithmically, with an
explicit differentiable remainder. This would also improve numerical
continuation near $\rho=1$.

\subsection{Divergent angular expansions}

\paragraph{13. Determine optimal truncation and the smallest attainable error.}
Theorem~\ref{angular:thm:allorders} keeps the cutoff fixed as
$a\to\infty$, whereas Theorem~\ref{angular:thm:divergence} rules out
ordinary convergence at fixed $a$. The next problem is quantitative:
how should the cutoff depend on $a$ to minimize the remainder, and
what is the resulting error scale? The proof must control the growth
of the finite implicit-function constants as the number of retained
modes increases. The prime-frequency obstruction supplies explicit
large coefficients and should constrain any proposed optimal rule.

\paragraph{14. Find a justified summation procedure.}
Can one recover the actual angular zero from the divergent coefficient
family by a natural summation method? Possible approaches include
keeping a finite Fourier equation exact while expanding only its low
modes, or reorganizing terms by a parameter introduced before the
mode cutoff. Any summability theorem must specify the order of
summation and explain why the unbounded grouped terms do not invalidate
the proposed procedure. A formal appeal to analyticity of the complete
Fourier series is insufficient.

\paragraph{15. Extend the angular theorem to other colors and depths.}
The current local root is governed by a dominant Fourier mode and
a nonzero real derivative, while the source's signed density proves
global uniqueness. Which depth-three or non-Gaussian color families
have both properties? Separate these two tasks: a local expansion can
be constructed before global uniqueness is known, but identifying it
with the intended geometric zero requires an independent global
argument. Multi-sign-change densities may produce several roots and
coupled singular limits, which would connect this direction with the
Lerch bifurcation analysis above.

\paragraph{A practical order of attack.}
The most immediate projects are the first residual levels in item~1,
constructive logarithmic evaluations in item~4, certified interior
events in item~9, and the first optimal-truncation estimates in item~13.
Each has a precise finite or analytic starting point supplied by this
article. The broader period-independence questions remain important,
but should be pursued with the formal-versus-numerical distinction
preserved at every stage.
